% ===== §11 =====
\section{The Shared Presupposition and Its Inversion}
\label{p24:sec:inversion}

\noindent
With the framework stated, the antithesis can do the work the map assigned it: to draw the traditions from their several coordinates toward the vacant one. It does this not by refuting each in turn, which would be tedious and would miss what unites them, but by exposing the single presupposition they share beneath their disagreements and inverting it. The six traditions dispute almost everything---whether culture is a whole or a web or a structure or an encoding of power, whether it means or reproduces or dominates---but underneath the dispute they agree, with the partial exception of the two that lean toward the generative, on one thing: that culture is, in its mode of being, a \emph{thing already made}. An acquirable heritage, a prior web of meaning, a deep structure operating behind the agents' backs, a congealed order of power into which subjects are inducted---each of these is a culture that stands \emph{before} the individual, who enters it, bears it, is formed or dominated by it, but does not, in the analysis, make it. The individual is the container of culture or its effect, never, at the moment the theory looks, its generating source.

\subsection*{The inversion}

The framework of the previous section inverts this presupposition at its root. Culture is not a prior object into which individuals are inducted; it is a process that relational subjects continually generate in their coupling. It is not a web the individual falls into but a weaving the relation performs, and---this is the decisive addition---a weaving that turns back to weave the relation in its turn, so that the subjects who generate the culture are themselves generated, as these subjects, by the culture they generate. The priority the six traditions grant to the made-thing over the making is exactly reversed: the making is primary, and the made-thing, the sedimented culture that the traditions take for the whole, is the residue the making leaves behind and the material the next round of making takes up. What the transmission tradition called an acquirable stock is the ash of a fire it never watched burn; what the interpretive tradition called a web already spun is a weaving caught at a late moment and mistaken for a finished textile; what the structural tradition called a deep structure behind the agents is the sedimented regularity of a generation the agents are even now carrying on.

This inversion is, in its logical form, the Hegelian moment the second line of the paper supplies, and naming it as such shows that the inversion is not a mere reversal of emphasis but a determinate movement. Culture is spirit exteriorized: the relation goes out of itself into a visible, sedimented form---a word, a rite, an institution---and this going-out is not the loss the traditions implicitly treat it as, a congealing of living process into dead thing, but the necessary path by which the relation comes to know itself. The exteriorized culture is the relation's self-objectification, in which it can see what it has become and from which it returns to itself altered, entering the next round of generation already changed by having beheld its own work. The sedimented forms the traditions describe are real, then, and the inversion does not deny them; it relocates them, as moments in a movement, from the status of the thing culture \emph{is} to the status of the residue culture \emph{leaves} and the mirror in which it knows itself. Exteriorization, not as fall but as self-knowledge: this is the form of the inversion, and it is why the paper can keep everything the traditions correctly saw while denying the priority they granted it.

\subsection*{What is kept and what is discarded}

The inversion is not a wholesale rejection, and fairness to the traditions requires that its takings and leavings be stated exactly. Each tradition saw something real; the framework keeps the seeing and discards only the presupposition of the made-thing's priority. The table records this, tradition by tradition.

\begin{center}
\small
\renewcommand{\arraystretch}{1.4}
\begin{tabularx}{\linewidth}{@{}>{\RaggedRight\arraybackslash}p{3.2cm} >{\RaggedRight\arraybackslash}X >{\RaggedRight\arraybackslash}X@{}}
\toprule
\textcolor{rosedeep}{\textbf{Tradition}} & \textcolor{rosedeep}{\textbf{Kept (absorbed)}} & \textcolor{rosedeep}{\textbf{Discarded / sublated}} \\
\midrule
Transmission \& collective representation & transmissibility and sedimentation are real features of culture & its \emph{static, already-made} mode of being \\
Symbol \& interpretation & culture is a web of meaning, relational, read by thick description & the absolutizing of the Symbolic; meaning as prior to the subjects \\
Structure \& reproduction & culture does reproduce and does have structure & the \emph{extractable / transferable} premise of cultural capital \\
Power as ontology & power is internal to all culture (the anti-naive lesson) & sublated: that culture \emph{is} Symbolic power-encoding with no outside \\
The dialogic (Bakhtin) & \emph{kept most}: dialogism, unfinishedness, generation in answering & only what it lacks is added---drive, the registers, the material ground \\
The Chinese \emph{wen--hua} & the verbal, generative force of \emph{hua}; process before product & its historical entanglement with a cultivation imposed from above \\
\bottomrule
\end{tabularx}
\end{center}

Read across, the table shows that the dialogic strand is the framework's near kin, from which most is taken, and the power tradition the one that is not discarded but \emph{sublated}---its clear-eyed insistence that power is internal to all culture is kept in full, while its conclusion that culture is therefore nothing but Symbolic power-encoding, with no outside, is negated by the relational Real that the framework installs in a register power cannot reach. The remaining three each surrender only the presupposition of the already-made, in the specific form each gave it: the frozen heritage, the pre-spun web, the extractable asset. With the shared presupposition inverted and each tradition's truth reassigned its place, the antithesis can proceed from this general reversal to the two specific developments the framework demands: the account of culture as the exteriorization of drive, and the dialectic of the material and the ideal on which the whole rests.
